\subsubsection*{04.04.05 Führung}
\phantomsection
\addcontentsline{toc}{subsubsection}{04.04.05 Führung}
\label{sec:04.04.05_Fuehrung}

%%% Situation 

\par Der \gls{wsr} war als Matrixorganisation mit internen Mitarbeitenden (\gls{ais}, \gls{fnz}) und sieben externen Zulieferern strukturiert. Ohne disziplinarische Befugnisse war die Zusammenarbeit vollständig durch laterale Führung geprägt. Die initiale Analysephase unterschätzte die tatsächliche Komplexität der gewachsenen \gls{r3}-Landschaft: Individualentwicklungen, Sonderlösungen und fehlendes Know-how führten zu erheblichen Risiken (Terminplanung, Ressourcen, Stabilität). Programmplanung musste aufgrund hoher \gls{ws}-Abhängigkeiten mehrfach angepasst werden. Welle 1 wurde in vier Unterwellen (1.0, 1.1, 1.2, 1.3) aufgeteilt, was Druck auf Koordination, Priorisierung und Parallelbetrieb erhöhte. Entscheidend waren nicht nur technische Lieferobjekte, sondern insbesondere Führung eines heterogenen, verteilten Projektumfelds, um Programmziele zu erreichen, Risiken zu beherrschen und Stabilität für Produktivsetzungen zu sichern.



